\documentclass[a4paper,12pt,fleqn]{article}
\usepackage[ngerman]{babel}
\usepackage[utf8]{inputenc}
\usepackage[T1]{fontenc}
\usepackage{amsmath,amssymb}
\usepackage{xcolor}
\usepackage{tikz}
\usetikzlibrary{automata,positioning,arrows.meta}

\tikzset{
    ->,>=Stealth,
    node distance=2.5cm,
    every state/.style={thick,minimum size=1cm},
    initial text={},
}

\newcommand{\changed}[1]{\textcolor{blue}{#1}}
\newcommand{\removed}[1]{\textcolor{red}{#1}}

\newcommand{\methodbox}[1]{%
\medskip
\noindent\fbox{%
    \begin{minipage}{0.94\linewidth}
    \textbf{Methodik / Definition.}\par\smallskip
    \small
    #1
    \end{minipage}%
}
\medskip\par
}

\newcommand{\examplebox}[1]{%
\medskip
\noindent\fbox{%
    \begin{minipage}{0.94\linewidth}
    \textbf{Beispiele.}\par\smallskip
    \small
    #1
    \end{minipage}%
}
\medskip\par
}

\newcommand{\taskbox}[1]{%
\medskip
\noindent\fbox{%
    \begin{minipage}{0.94\linewidth}
    \textbf{Aufgabenstellung.}\par\smallskip
    \small
    #1
    \end{minipage}%
}
\medskip\par
}

\begin{document}
\raggedright
\setlength{\parindent}{0pt}
\setlength{\mathindent}{0pt}

{\LARGE\bfseries ETI -- Aufgabenblatt 5 \\ Praesenzaufgaben\par}
\vspace{2em}

\section*{Aufgabe 5.1}

\taskbox{
Gegeben ist die Grammatik $G_1=(V,\Sigma,R,S)$ mit
$V=\{S,A,B,X\}$, $\Sigma=\{0,1\}$ und
\[
\begin{aligned}
S &\to ABA\mid ABB\mid XA\mid B1,\\
B &\to 0\mid 1\mid A,\\
A &\to 0\mid 1\mid \varepsilon.
\end{aligned}
\]
Entfernen Sie alle $\varepsilon$-Regeln von nicht-Start-Nonterminalen.
}

\methodbox{
Zuerst bestimmen wir die nullable Nonterminale. Fuer jede Produktion fuegen wir dann alle
Varianten hinzu, in denen nullable Vorkommen gestrichen werden. Die Regel
$S\to\varepsilon$ darf bleiben, weil $S$ das Startsymbol ist und die urspruengliche Grammatik
das leere Wort erzeugen kann.
}

Nullable sind
\[
    A,\qquad B \quad\text{und}\quad S.
\]
Dabei gilt $A\Rightarrow\varepsilon$, $B\Rightarrow A\Rightarrow\varepsilon$ und
$S\Rightarrow ABA\Rightarrow\varepsilon$.

Die Produktionen nach dem Entfernen der $\varepsilon$-Regel $A\to\varepsilon$ lauten:

\begin{flushleft}
\scriptsize
\begin{tabular}{@{}l@{\hspace{1.2em}}l@{}}
\textbf{Vorher} & \textbf{Nachher ohne $\varepsilon$-Regeln}\\[0.4em]
$\begin{aligned}
S &\to ABA\mid ABB\mid XA\mid B1\\
B &\to 0\mid 1\mid A\\
A &\to 0\mid 1\mid \removed{\varepsilon}
\end{aligned}$
&
$\begin{aligned}
S &\to ABA \mid \changed{BA} \mid \changed{AB} \mid \changed{AA}
      \mid \changed{B} \mid \changed{A} \mid ABB \mid \changed{BB}\\
  &\mid XA \mid \changed{X} \mid B1 \mid \changed{1} \mid \changed{\varepsilon}\\
B &\to 0 \mid 1 \mid A\\
A &\to 0 \mid 1
\end{aligned}$
\end{tabular}
\end{flushleft}
Damit gibt es keine $\varepsilon$-Regel eines Nicht-Startsymbols mehr.

\section*{Aufgabe 5.2}

\taskbox{
Gegeben ist die Grammatik $G_2=(V,\Sigma,R,S)$ mit
$V=\{S,T,U,V,W\}$, $\Sigma=\{k,j,h,b\}$ und
\[
\begin{aligned}
S&\to T\mid Vj,\\
T&\to U\mid b\mid Wh,\\
U&\to W\mid k,\\
V&\to W\mid j\mid Vb,\\
W&\to h.
\end{aligned}
\]
Entfernen Sie alle Kettenregeln.
}

\methodbox{
Eine Kettenregel hat die Form $A\to B$ mit $A,B\in V$. Fuer jedes Nonterminal $A$ betrachten
wir alle Nonterminale, die von $A$ nur mit Kettenregeln erreichbar sind, und uebernehmen deren
nicht-Kettenregeln direkt nach $A$.
}

Die Kettenabschluesse sind
\[
\begin{array}{c|l}
A & \text{ueber Kettenregeln erreichbar}\\ \hline
S & S,T,U,W\\
T & T,U,W\\
U & U,W\\
V & V,W\\
W & W
\end{array}
\]

Nach dem Entfernen aller Kettenregeln erhalten wir:

\begin{flushleft}
\scriptsize
\begin{tabular}{@{}l@{\hspace{1.2em}}l@{}}
\textbf{Vorher} & \textbf{Nachher ohne Kettenregeln}\\[0.4em]
$\begin{aligned}
S&\to \removed{T}\mid Vj\\
T&\to \removed{U}\mid b\mid Wh\\
U&\to \removed{W}\mid k\\
V&\to \removed{W}\mid j\mid Vb\\
W&\to h
\end{aligned}$
&
$\begin{aligned}
S &\to Vj \mid \changed{b} \mid \changed{Wh} \mid \changed{k} \mid \changed{h}\\
T &\to b \mid Wh \mid \changed{k} \mid \changed{h}\\
U &\to k \mid \changed{h}\\
V &\to j \mid Vb \mid \changed{h}\\
W &\to h
\end{aligned}$
\end{tabular}
\end{flushleft}

\section*{Aufgabe 5.3}

\taskbox{
Gegeben ist die Grammatik $G_3=(V,\Sigma,R,S)$ mit
$V=\{S,X,Y,Z\}$, $\Sigma=\{a,b,c,d\}$ und
\[
\begin{aligned}
S &\to XYZ\mid aX,\\
X &\to XYZ\mid a,\\
Y &\to YYZ\mid Zb,\\
Z &\to a\mid ab.
\end{aligned}
\]
Bringen Sie die Grammatik in Chomsky-Normalform.
}

\methodbox{
In Chomsky-Normalform haben alle Regeln die Form $A\to BC$ oder $A\to a$.
Terminale in rechten Seiten der Laenge mindestens $2$ ersetzen wir daher durch eigene
Variablen. Rechte Seiten der Laenge $3$ werden binaer zerlegt.
}

Es gibt keine $\varepsilon$-Regeln und keine Kettenregeln. Wir fuehren
\[
    T_a\to a,\qquad T_b\to b,\qquad C_{YZ}\to YZ
\]
ein. Die rechten Seiten $XYZ$ und $YYZ$ koennen dann als $XC_{YZ}$ bzw. $YC_{YZ}$
geschrieben werden.

\medskip
Eine Chomsky-Normalform ist also
\[
G_{\mathrm{CNF}}=(V',\{a,b,c,d\},R',S)
\]
mit
\[
V'=\{S,X,Y,Z,T_a,T_b,C_{YZ}\}
\]
und Regeln

\begin{flushleft}
\scriptsize
\begin{tabular}{@{}l@{\hspace{1.2em}}l@{}}
\textbf{Vorher} & \textbf{Nachher in CNF}\\[0.4em]
$\begin{aligned}
S &\to \removed{XYZ}\mid \removed{aX}\\
X &\to \removed{XYZ}\mid a\\
Y &\to \removed{YYZ}\mid \removed{Zb}\\
Z &\to a\mid \removed{ab}
\end{aligned}$
&
$\begin{aligned}
S &\to \changed{XC_{YZ}}\mid \changed{T_aX}\\
X &\to \changed{XC_{YZ}}\mid a\\
Y &\to \changed{YC_{YZ}}\mid \changed{ZT_b}\\
Z &\to a\mid \changed{T_aT_b}\\
C_{YZ} &\to \changed{YZ}\\
T_a &\to \changed{a}\\
T_b &\to \changed{b}
\end{aligned}$
\end{tabular}
\end{flushleft}
Alle Produktionen haben nun die Form $A\to BC$ oder $A\to a$.

\section*{Aufgabe 5.4}

\taskbox{
Betrachten Sie
\[
M=\{w\in\{0,1\}^*\mid |w|_0=|w|_1\}.
\]
Konstruieren Sie einen PDA $P$ mit $L(P)=M$.
}

\methodbox{
Der Stack speichert nur die noch nicht ausgeglichenen Zeichen. Liest der PDA ein Zeichen,
das zum aktuellen Stacktop passt, wird es gepusht. Liest er das jeweils andere Zeichen,
wird das Stacktop geloescht. Am Ende ist genau dann nur noch das Bodensymbol uebrig,
wenn gleich viele $0$en und $1$en gelesen wurden.
}

Wir verwenden das Stackalphabet
\[
    \Gamma=\{0,1,\$\}
\]
mit Bodensymbol $\$$ und akzeptieren per Endzustand. Der PDA ist
\[
P=(\{q_0,q_1,q_2\},\{0,1\},\{0,1,\$\},\delta,q_0,\{q_2\}).
\]

\begin{center}
\begin{tikzpicture}[node distance=4.1cm]
    \node[state, initial] (q0) {$q_0$};
    \node[state, right=of q0] (q1) {$q_1$};
    \node[state, accepting, right=of q1] (q2) {$q_2$};

    \path (q0) edge node[above] {$\varepsilon,\varepsilon\to\$$} (q1)
          (q1) edge[loop above] node {$\begin{array}{c}
              0,1\to\varepsilon\quad 1,0\to\varepsilon
          \end{array}$} ()
          (q1) edge[loop below] node {$\begin{array}{c}
              0,\$\to 0\$\quad 1,\$\to 1\$\\
              0,0\to 00\quad 1,1\to 11
          \end{array}$} ()
          (q1) edge node[above] {$\varepsilon,\$\to\varepsilon$} (q2);
\end{tikzpicture}
\end{center}

\examplebox{
Fuer $w=0011$ gleicht sich alles aus:
\[
    \$ \xrightarrow{0} 0\$ \xrightarrow{0} 00\$
    \xrightarrow{1} 0\$ \xrightarrow{1} \$.
\]
Am Ende liegt nur noch das Bodensymbol $\$$ auf dem Stack, also wird akzeptiert.

Auch bei gemischter Reihenfolge ist nur die Anzahl wichtig. Fuer $w=0101$ gilt:
\[
    \$ \xrightarrow{0} 0\$ \xrightarrow{1} \$
    \xrightarrow{0} 0\$ \xrightarrow{1} \$.
\]
Also wird auch $0101$ akzeptiert.

Fuer $w=001$ bleibt dagegen eine $0$ uebrig:
\[
    \$ \xrightarrow{0} 0\$ \xrightarrow{0} 00\$
    \xrightarrow{1} 0\$.
\]
Der Stack ist nicht leer bis auf $\$$, also wird nicht akzeptiert.

Ein gemischtes nicht akzeptiertes Beispiel ist $w=011$:
\[
    \$ \xrightarrow{0} 0\$ \xrightarrow{1} \$
    \xrightarrow{1} 1\$.
\]
Hier bleibt eine $1$ uebrig, also wird nicht akzeptiert.

Noch ein gemischtes Beispiel ist $w=10100$:
\[
    \$ \xrightarrow{1} 1\$ \xrightarrow{0} \$
    \xrightarrow{1} 1\$ \xrightarrow{0} \$
    \xrightarrow{0} 0\$.
\]
Hier bleibt eine $0$ uebrig, also wird ebenfalls nicht akzeptiert.
}

Nach jedem gelesenen Praefix enthaelt der Stack ueber $\$$ genau die Differenz der bisher
gelesenen Zeichen: lauter $0$en, falls bisher mehr $0$en gelesen wurden, oder lauter $1$en,
falls bisher mehr $1$en gelesen wurden. Deshalb ist nach dem ganzen Wort genau dann nur noch
$\$$ auf dem Stack, wenn $|w|_0=|w|_1$ gilt. Genau dann kann der PDA nach $q_2$ wechseln.

\section*{Aufgabe 5.5}

\taskbox{
Vervollstaendigen Sie eine kontextsensitive Grammatik fuer
\[
    A=\{a^n b^n c^n\mid n\in\mathbb{Z}_{\ge 0}\}
\]
ausgehend von
\[
S\to A\mid \varepsilon,\qquad
A\to abc\mid aABc.
\]
}

\methodbox{
Die Regel $A\to aABc$ erzeugt pro Rekursionsschritt ein weiteres $a$, ein weiteres $B$
und ein weiteres $c$. Danach muessen die $B$-Symbole nach links ueber die $c$-Symbole
wandern und dort zu $b$ werden.
}

Eine passende Vervollstaendigung ist:
\[
\begin{aligned}
S &\to A\mid \varepsilon,\\
A &\to abc\mid aABc,\\
cB &\to Bc,\\
bB &\to bb.
\end{aligned}
\]

Alle Regeln sind kontextsensitiv: Sie verkuerzen nicht. Die einzige $\varepsilon$-Regel ist
$S\to\varepsilon$, und $S$ kommt auf keiner rechten Seite vor.

\medskip
Nach $n-1$ Anwendungen von $A\to aABc$ und anschliessender Anwendung von $A\to abc$
entsteht fuer $n\ge 1$
\[
    a^n b\,c(Bc)^{n-1}.
\]
Mit $cB\to Bc$ werden alle $B$-Symbole nach links vor die $c$-Symbole geschoben:
\[
    a^n b\,c(Bc)^{n-1}
    \Rightarrow^*
    a^n bB^{n-1}c^n.
\]
Danach wandelt $bB\to bb$ die $B$-Symbole von links nach rechts in $b$ um:
\[
    a^n bB^{n-1}c^n
    \Rightarrow^*
    a^n b^n c^n.
\]
Fuer $n=0$ erzeugt $S\to\varepsilon$ das leere Wort. Somit erzeugt die Grammatik genau
\[
    \{a^n b^n c^n\mid n\in\mathbb{Z}_{\ge 0}\}.
\]

\end{document}
